\documentclass[12pt,a4paper]{article}
\PassOptionsToPackage{table}{xcolor}

% ----- GROUPED PACKAGES -----
% General Utilities
\usepackage[utf8]{inputenc}
\usepackage{geometry}
\usepackage{float}
\usepackage{afterpage}
\usepackage{enumitem}
\usepackage{ragged2e}
\usepackage[normalem]{ulem}
\usepackage{hyperref}
\hypersetup{colorlinks=true, urlcolor=blue, allcolors=blue}
% Graphics and Drawing
\usepackage{graphicx}
\usepackage{pgf}
\usepackage{tikz}
\usepackage{svg}
\usepackage{stanli}
\usepackage{wrapfig}
\usepackage{rotating}
\usepackage[table]{xcolor}

% Math and Science
\usepackage{amsmath}
\usepackage{siunitx}

% Tables and Figures
\usepackage{booktabs}
\usepackage{caption}
\usepackage{subcaption}
\usepackage{multirow}

% Document Structure and Layout
\usepackage{titlesec}
\usepackage{titling}
\usepackage{pgfpages}
\usepackage{changepage}

% Code Listings
\usepackage{listings}


% ----- CUSTOM DEFINITIONS -----

% Custom list settings
\setlist[itemize]{noitemsep, topsep=0pt, parsep=0pt, partopsep=0pt, leftmargin=*,
                 labelsep=0.5em, itemsep=1pt, label=\textbullet, font=\small}

% Custom section spacing
\titlespacing*{\section}{2pt}{*2}{*2}
\titlespacing*{\subsection}{3pt}{*3}{*3}

% Custom table row colors
\definecolor{tablerow1}{RGB}{240,240,240}
\definecolor{tablerow2}{RGB}{220,230,240}

% Enhanced VS Code Inspired with Light Background
\definecolor{vscode-blue}{RGB}{0,0,255}
\definecolor{vscode-green}{RGB}{0,128,0}
\definecolor{vscode-red}{RGB}{163,21,21}
\definecolor{vscode-purple}{RGB}{163,53,238}
\definecolor{vscode-orange}{RGB}{200,90,30}

\lstdefinestyle{vscodelight}{
    basicstyle=\ttfamily\footnotesize\bfseries\color[RGB]{30,30,30},
    backgroundcolor=\color[RGB]{240,248,255},
    rulecolor=\color[RGB]{200,200,200},
    frame=single,
    frameround=tttt,
    framesep=6pt,
    xleftmargin=13pt,
    xrightmargin=13pt,
    numbers=left,
    numberstyle=\tiny\bfseries\color[RGB]{100,100,100},
    numbersep=10pt,
    breaklines=true,
    breakatwhitespace=true,
    showstringspaces=false,
    tabsize=6,
    captionpos=b,
    keywordstyle=\bfseries\color{vscode-blue},
    commentstyle=\itshape\bfseries\color{vscode-green},
    stringstyle=\bfseries\color{vscode-red},
    identifierstyle=\bfseries\color[RGB]{30,30,30},
    emphstyle=\bfseries\color{vscode-purple},
    emph={function,class,return,import,from,as},
    moredelim=[is][\bfseries\color{vscode-orange}]{@}{@},
    linewidth=1.1\linewidth,
    xleftmargin=0pt,
    xrightmargin=0pt
}

\lstset{style=vscodelight}

% Additional formatting for academic papers
\lstset{
    aboveskip=0.7\baselineskip,
    belowskip=0.7\baselineskip,
    abovecaptionskip=5pt,
    belowcaptionskip=5pt,
    escapechar=!,
    prebreak=\raisebox{0ex}[0ex][0ex]{\ensuremath{\hookleftarrow}},
    postbreak=\raisebox{0ex}[0ex][0ex]{\ensuremath{\hookrightarrow}},
    extendedchars=true
}

% For even wider code blocks
\lstset{
    frame=lrtb,
    framerule=0.5pt
}

% Set page margins
\geometry{a4paper, margin=1in}

% Customize title style
\pretitle{\begin{center}\LARGE\bfseries}
\posttitle{\par\end{center}\vskip 0.3em}
\preauthor{\begin{center}\large}
\postauthor{\end{center}}
\predate{\begin{center}\large}
\postdate{\end{center}}


\begin{document}

% Title page
\begin{titlepage}
    \centering
    \vspace{1cm}

    \includegraphics[width=0.25\linewidth]{university.png}\\
    \vspace{0.4cm}
    \includegraphics[width=0.2\linewidth]{entc1.png}\\
    
    \vspace{0.4cm}
    {\Large University of Moratuwa \par}
     \vspace{0.2cm}
     {\Large Department of Electronic and Telecommunication \par}

    \vspace{1cm}
    
    \rule{\textwidth}{0.4pt}
    \vspace{1cm}
    
    {\LARGE\bfseries Assignment 01 \par}
    \vspace{1cm}
    {\Large EN4554 - Deep Learning for Vision \par}
    \vspace{2cm}
    \rule{\textwidth}{0.4pt}
    \begin{center}
    \vspace{2cm}
    \textbf{\Large Sulojan.R - 220626V}\\
    \vspace{0.5cm}
    \textbf{\Large Praveen V.V.J. - 220491B}\\
\end{center}
\vfill


\vfill
    \begin{center}
        {\Large 4th of November 2026}
    \end{center}
\end{titlepage}

% ===================== Part 1 (Sulojan R.) goes here =====================


% ===================== Part 2: LLM Agents, Tool Calling ====================
\newgeometry{top=0.5in, bottom=0.55in, left=0.6in, right=0.6in}
\begingroup
\titleformat*{\section}{\large\bfseries}
\titleformat*{\subsection}{\normalsize\bfseries}
\titlespacing*{\section}{0pt}{0pt}{4pt}
\titlespacing*{\subsection}{0pt}{6pt}{2pt}
\setlength{\topsep}{2pt}
\section*{Part 2: LoRA Fine-tuning a Small LLM for Tool Calling}
\footnotesize
\setlength{\parskip}{2pt}
\setlength{\parindent}{0pt}

\subsection*{1. Setup}
\begin{itemize}
    \item \textbf{Model:} Qwen2.5-0.5B-Instruct, used through its built-in tool-calling chat template (each call is a JSON object wrapped in \verb|<tool_call>| tags). Greedy decoding, at most 128 new tokens per turn.
    \item \textbf{Tools (4):} \verb|calculator| (arithmetic), \verb|unit_convert| (length, mass, $^\circ$C/$^\circ$F), \verb|get_weather| (made-up weather for six Sri Lankan cities) and \verb|lookup_order| (60 made-up orders). The invented values force the model to call the tools instead of answering from memory.
    \item \textbf{Agent loop:} generate $\rightarrow$ parse any tool call $\rightarrow$ execute it $\rightarrow$ append the result as a tool message $\rightarrow$ repeat, until a plain-text reply or 4 steps.
    \item \textbf{Data:} template-generated, with gold answers computed by running the tools. \textbf{Train: 1{,}200} requests (856 single-call, 118 two-call chains, 226 needing no tool). \textbf{Test: 100} requests (68 single-call, 5 chains, 27 no-tool; 77 unique), built from separately written templates with different wording. No test request appears in the training set.
    \item \textbf{LoRA:} $r=16$, $\alpha=32$, dropout 0.05 on the q, k, v, o, gate, up and down projections (8.80\,M trainable parameters, 1.75\,\%). Loss on assistant tokens only; system, user and tool-result tokens are masked. 2 epochs, learning rate $2\times10^{-4}$ (cosine, 5 warm-up steps), effective batch 16 (micro-batch $2\times8$ accumulation, because batch 8 ran out of memory with Qwen's 152k-token vocabulary), fp16 mixed precision, seed 42. 150 steps in 6\,min\,42\,s on a Colab T4; training loss fell from 0.34 to 0.0003.
\end{itemize}

\subsection*{2. Baseline vs.\ LoRA (same 100-request test set)}
\begin{center}
\begin{tabular}{@{}llrrr@{}}
    \toprule
    \textbf{Metric} & \textbf{Definition (denominator)} & \textbf{Baseline} & \textbf{LoRA} & \textbf{Change} \\
    \midrule
    Call validity & made $\geq 1$ call, all well-formed (tool requests, $n=73$) & 83.6\,\% & \textbf{100.0\,\%} & +16.4 \\
    Call accuracy & right tool(s) and right arguments ($n=73$) & 68.5\,\% & \textbf{90.4\,\%} & +21.9 \\
    Abstention & answered without a tool when none was needed ($n=27$) & 59.3\,\% & \textbf{96.3\,\%} & +37.0 \\
    Task success & final answer contains the gold facts, end to end ($n=100$) & 55.0\,\% & \textbf{92.0\,\%} & +37.0 \\
    \bottomrule
\end{tabular}
\end{center}
The baseline never produced malformed JSON: all of its validity failures were tool requests where it made no call at all.

\subsection*{3. Failure analysis}
\noindent
\begin{minipage}[t]{0.37\linewidth}
    \vspace{0pt}
    \begin{tabular}{@{}lrr@{}}
        \toprule
        \textbf{Category} & \textbf{Base} & \textbf{LoRA} \\
        \midrule
        Ignored tool result & 15 & 0 \\
        Missed call & 14 & 0 \\
        Unnecessary call & 11 & 1 \\
        Wrong no-tool answer & 7 & 0 \\
        Wrong tool & 6 & 4 \\
        Wrong arguments & 3 & 3 \\
        Wrong format & 0 & 0 \\
        \midrule
        \textbf{Total failures} & \textbf{56} & \textbf{8} \\
        \bottomrule
    \end{tabular}

    \vspace{3pt}
    One category per request; ``missed call'' includes chains cut short. On the same requests, LoRA \textbf{fixed 49} of the 56 baseline failures, \textbf{left 7 unfixed} and \textbf{introduced 1}. All 8 LoRA failures were inspected (only 8 exist, fewer than the 10 asked for), along with 10 baseline failures.
\end{minipage}\hfill
\begin{minipage}[t]{0.60\linewidth}
    \vspace{0pt}
    \textbf{Baseline errors were decision errors, not format errors.} It asked the user for data instead of calling (``should I pack a coat?'' $\rightarrow$ ``could you provide the current temperature?''), called tools for general knowledge (\verb|lookup_order(order_id="RAM")| for ``What does RAM stand for?''), refused easy questions (``I can only provide details about mathematical operations and weather'') and misreported correct results (``the backpack is expected to be back in stock in 6 days''). \textbf{LoRA fixed} every ignored result, refusal and wrong-tool error, 10 of the 11 unnecessary calls and 10 of the 14 missed calls.

    \textbf{Not fixed (7):} ``Take $b$ away from $a$'' has wrong arguments in both models, and LoRA computes $b-a$ (2 cases). The boiling-point question still triggers \verb|calculator("0*100")|, answered as ``0 $^\circ$C'' (1). ``I'm travelling to \textit{city}, should I pack a coat?'' moved from a missed call to the wrong tool, \verb|calculator("1*24")|, answered as ``24 hours until you arrive'' (4). \textbf{Introduced (1):} ``Vavuniya weather report, please'' now calls \verb|get_weather(city="Lisbon")|.
\end{minipage}

\vspace{4pt}
The remaining LoRA errors share one cause: \textbf{over-fitting to the training templates} (loss $\approx 2\times10^{-4}$ by epoch 2). The model keys on surface patterns, such as reading ``travelling to \textit{city}'' as the trained ``how many hours until \ldots{} arrives'' chain, and operand order in subtraction, rather than on intent. More varied phrasings, paraphrase augmentation or a single epoch would be the next step. The ignored-result count is also somewhat inflated by a strict answer check that requires both the order status and the item.

\subsection*{4. Would a better prompt alone have been enough?}
Only partly: the baseline already produced well-formed calls through Qwen's built-in tool template, so its failures were decisions (whether to call a tool at all in 32 cases, how to report the result in 15), and a stricter system prompt with a few worked examples could plausibly reduce these without any training. However, a 0.5\,B model follows long instructions unreliably, and LoRA removed 49 of the 56 failures after under seven minutes of training (we did not run a prompt-only ablation to measure the gap). The errors LoRA left come from template over-fitting, which better training data, not a better prompt, would address.

\endgroup
\restoregeometry

\end{document}